\documentclass[conference]{IEEEtran}
\IEEEoverridecommandlockouts
% If you don't need a funding footnote, you can comment out the line above.

\usepackage{cite}
\usepackage{amsmath,amssymb,amsfonts}
\usepackage{graphicx}
\usepackage{textcomp}

\def\BibTeX{{\rm B\kern-.05em{\sc i\kern-.025em b}\kern-.08em
    T\kern-.1667em\lower.7ex\hbox{E}\kern-.125emX}}

\begin{document}

\title{Detecting Bot-Generated Comments: Comparing Transformer-Based and Traditional Machine Learning Detectors}

\author{\IEEEauthorblockN{James Kelly}
\IEEEauthorblockA{\textit{Department of Computer Science and Information Systems}\\
\textit{University of Limerick}\\
Limerick, Ireland \\
23370858}
}

\maketitle

\begin{abstract}
    Online platforms increasingly host comments generated by bots or large language models alongside genuine human opinions. This mixture creates challenges for platform trust, moderation, and users' ability to judge authenticity. This paper examines the problem of distinguishing bot-generated comments from human comments on online platforms and compares traditional machine learning methods with transformer-based detectors in a supervised classification setting. By evaluating both model families under the same conditions, the study aims to identify which approach is more effective for realistic comment-level bot detection.
\end{abstract}

\section{Introduction}

Online platforms are increasingly flooded with comments that are not written by humans but generated by bots or large language models.
What once were simple discussion threads now mix genuine opinions with automated praise, spam, and subtle persuasion.
This growing presence of machine-generated comments raises concerns for platform trust, moderation workloads, and users’ ability to judge what is authentic \cite{zaitsu_can_2024,li_mage_2024}.
Before any detector can be evaluated on accuracy or robustness, we need to understand how well different model families can separate bot-generated comments from human ones in the first place \cite{kumar_m_scalar_2024,ahmad_human_2026}.

Recent literature shows that machine-generated text detection is now an active area of research, with approaches ranging from classical machine learning and stylometric analysis to transformer-based and ensemble models. 
Kumar et al. demonstrate that a RoBERTa ensemble can perform strongly in a shared-task setting, while Li et al. show that detection becomes more difficult in broader, 
more realistic settings where texts come from multiple domains and generation sources \cite{kumar_m_scalar_2024,li_mage_2024}. 
At the comment level, Zaitsu et al. show that fake public comments generated by ChatGPT can be distinguished from human comments using stylometric features and classification methods \cite{zaitsu_can_2024}. 
Ahmad et al.'s recent survey also indicates that transformer-based detectors are currently among the strongest approaches in the wider machine-generated text detection literature \cite{ahmad_human_2026}.
Recent work also shows that generative AI can alter the quality, authenticity, and volume of social-media content, which reinforces the need to study detection in realistic platform settings \cite{moller_impact_2026}.
This strengthens the case for evaluating detectors in a comment-level setting, where short, informal, and context-dependent language makes the task more challenging than general machine-generated text detection.

However, the existing literature does not yet provide a clear comparison of traditional machine learning methods and transformer-based detectors specifically for bot-generated comments in realistic online-platform contexts.
Much of the current work either focuses on general machine-generated text detection or on domain-specific datasets that may not fully reflect the short, informal, and varied nature of online comments. As a result, there is still a need for a focused comparison that tests whether traditional feature-based approaches or transformer-based models are more robust when applied to bot-generated comments. This study contributes a focused comment-level comparison of traditional machine learning and transformer-based detectors, with the goal of identifying which approach is more effective in realistic online-platform settings. The study will frame this comparison as a supervised classification task using labelled bot-generated and human-written comments, allowing the two model families to be evaluated under the same conditions.

The study asks whether traditional feature-based classifiers or transformer-based models are more effective for comment-level bot detection, and it addresses this question through a supervised comparative evaluation of labelled bot-generated and human-written comments.
The remainder of this paper is structured as follows: the next section outlines the literature context and research gap, followed by the proposed comparative approach and the intended evaluation strategy.
\bibliographystyle{IEEEtran}
\bibliography{345}

\end{document}